% page 292 of the facsimile (image p-291 of the scan)
% block 157.5 x 229.0 mm ; left margin 8.0 mm
\pgc{-12.700mm}{0.000mm}{
\l{\cel{3}284}
\l{}
\l{}
\l{\sou{kokoso}. (\sou{bot.})\cel{2}Frukto di granda arboro ek la familio \textquotedbl{}pal-}
\l{\cel{3}mieri\textquotedbl{}. L.\cel{1}\sou{cocus nucifera}. - DEFIRS.}
\l{}
\l{\sou{kolo.}\cel{3}Parto di la korpo di la homo, di animalo, qua unio-}
\l{\cel{3}nas la kapo a la torso. - FIS.}
\l{}
\l{\sou{kolao}.\cel{2}(\sou{bot.)}\cel{2}Arboro de la regioni tropikala di Afrika, qua}
\l{\cel{3}produktas \textquotedbl{}nuco\textquotedbl{} uzata kom tonizivo por la kordio e kom}
\l{\cel{3}stimulivo nervala. - L.\cel{1}\sou{cola acuminata}. - DEFI.}
\l{}
\l{\sou{kolacionar}.\cel{2}(\sou{trans., kun}). Komparar ica kun ita,o kun la ori-}
\l{\cel{3}ginalo, kopiurin, reprodukturin de texto, kom verifiko}
\l{\cel{3}di lia exakteso. - DEFI.}
\l{}
\l{+kolapsar. (netrans.)\cel{2}(aludante homo) . Falar pro exhausteso,}
\l{\cel{3}extenueso, febleso. - E.}
\l{}
\l{\sou{kolareto.}\cel{2}Stofo-garnituro por la kolo, generale plisoza, quan}
\l{\cel{3}+weris la viri lor la rejo Henriko IV di Francia (t. e.}
\l{\cel{3}cirkum la yaro 1600) e la mulieri, divers-epoke. - DEFI.}
\l{}
\l{\sou{kolaterala}.\cel{2}I. Situita laterale, relateulo. - II. (\sou{metaf.})}
\l{\cel{3}Qua apartenas a la genealogio-lineo \textquotedbl{}frati, kuzi, nevi,}
\l{\cel{3}onkli\textquotedbl{}. - DEFIS.}
\l{}
\l{\sou{kolchiko}. (\sou{bot.)}\cel{3}Planto bulboza, kun flori violea, di qua}
\l{\cel{3}la toxikeso uzesas por kuracar artritismo. - L.\cel{1}\sou{colchicum}}
\l{\cel{3}\sou{autumnale}. - EFIS.}
\l{}
\l{\sou{kolda.}\cel{2}I. Qua kontenas poka kaloro. - II. (\sou{metaf.})\cel{2}Qua ne}
\l{\cel{3}esas ardoroza, nek pasionoza. - DE.}
\l{}
\l{\sou{koldkremo}.\cel{2}Kosmetiko por la pelo, qua konsistas ye vaxo}
\l{\cel{3}blanka, spermaceto, ed oleo ek mandeli. - DEFIRS.}
\l{}
\l{\sou{kolego}. - Singla de la homi qui havas ofico publika, civila}
\l{\cel{3}od armeala, egardata relate a la ceteri. - DEFIRS.}
\l{}
\l{\sou{kolegio.}\cel{2}I. Korpo ek personi submisata a regularo komuna :}
\l{\cel{3}1. (en\cel{1}\sou{Roma, antique})\cel{2}Korporaciono ; 2. la korpo ek}
\l{\cel{3}elekteri politikala qui apartenas ad un distrikto votala;}
\l{\cel{3}3. edukerio publika - en Francia : edukerio municipala}
\l{\cel{3}(kontraste kun la \textquotedbl{}liceo\textquotedbl{}, edukerio statala). - DEFIRS.}
\l{}
\l{\sou{kolektar}.\cel{2}(\sou{trans.}) Asemblar (por facar ek li ensemblo) kozi}
\l{\cel{3}amasigita ek fonti diversa. - DEFIRS.}
\l{}
\l{\sou{kolektivismo}. (\sou{sociologio)}.\cel{2}Sistemo qua ofras, kom solvuro}
\l{\cel{3}di la problemo \textquotedbl{}quale desaparigar povreso\textquotedbl{}, la komunigo,}
\l{\cel{3}profite da socio, di omna produktili. - DEFIRS.}
}
